\[ %%%%%%%%%%%%%%%%%%%%%%%%%%%%%%%%%%%%%%%%%%%%%%%%%%%%%%%%%%%%%%%%%%%%%%%%%%%%%%%% \subsection{Dataset Generation} \label{sec:dataset_generation} %%%%%%%%%%%%%%%%%%%%%%%%%%%%%%%%%%%%%%%%%%%%%%%%%%%%%%%%%%%%%%%%%%%%%%%%%%%%%%%% The evaluation of PHYSGEN requires datasets that preserve the underlying physical structure of nonlinear dynamical systems. Instead of treating observations as independent samples, PHYSGEN considers data as continuous physical trajectories generated by an underlying evolution process. The general dataset formulation is: \begin{equation} \mathcal{D} = \{ G_0,G_1,\dots,G_T \}, \end{equation} where each graph state $G_t$ represents the physical configuration at time $t$. %%%%%%%%%%%%%%%%%%%%%%%%%%%%%%%%%%%%%%%%%%%%%%%%%%%%%%%%%%%%%%%%%%%%%%%%%%%%%%%% \subsubsection{Physical Trajectory Generation} %%%%%%%%%%%%%%%%%%%%%%%%%%%%%%%%%%%%%%%%%%%%%%%%%%%%%%%%%%%%%%%%%%%%%%%%%%%%%%%% For each benchmark system, trajectories are generated by solving the governing equations using numerical integration methods. The continuous physical evolution: \begin{equation} \frac{dX}{dt} = F(X,t,\Theta) \end{equation} is transformed into discrete observations: \begin{equation} X_{t+1} = X_t + \Delta t F(X_t,\Theta). \end{equation} The resulting trajectory is: \begin{equation} \mathcal{T} = \{X_0,X_1,\dots,X_T\}. \end{equation} %%%%%%%%%%%%%%%%%%%%%%%%%%%%%%%%%%%%%%%%%%%%%%%%%%%%%%%%%%%%%%%%%%%%%%%%%%%%%%%% \subsubsection{Lorenz-63 Dataset Generation} %%%%%%%%%%%%%%%%%%%%%%%%%%%%%%%%%%%%%%%%%%%%%%%%%%%%%%%%%%%%%%%%%%%%%%%%%%%%%%%% For the Lorenz-63 system, trajectories are generated using numerical integration of the nonlinear ODE system. The state vector is: \begin{equation} X_t=[x_t,y_t,z_t]^T . \end{equation} A fourth-order Runge--Kutta method (RK4) is used: \begin{equation} X_{t+1} = RK4(X_t,\Delta t). \end{equation} Multiple trajectories are generated using different initial conditions: \begin{equation} X_0^{(i)} \sim \mathcal{P}_{init}. \end{equation} This provides diverse examples of chaotic evolution. %%%%%%%%%%%%%%%%%%%%%%%%%%%%%%%%%%%%%%%%%%%%%%%%%%%%%%%%%%%%%%%%%%%%%%%%%%%%%%%% \subsubsection{Reaction-Diffusion Dataset Generation} %%%%%%%%%%%%%%%%%%%%%%%%%%%%%%%%%%%%%%%%%%%%%%%%%%%%%%%%%%%%%%%%%%%%%%%%%%%%%%%% For spatial dynamical systems, the continuous field is discretized into a finite computational grid. The field representation is: \begin{equation} U_t = [u_1,u_2,\dots,u_N]. \end{equation} Each spatial location becomes a graph node: \begin{equation} v_i \in V. \end{equation} Neighboring grid elements define graph connectivity: \begin{equation} (i,j)\in E \end{equation} when two spatial points interact. The generated dataset is therefore: \begin{equation} G_t = (V,E,X_t). \end{equation} %%%%%%%%%%%%%%%%%%%%%%%%%%%%%%%%%%%%%%%%%%%%%%%%%%%%%%%%%%%%%%%%%%%%%%%%%%%%%%%% \subsubsection{Navier--Stokes Dataset Generation} %%%%%%%%%%%%%%%%%%%%%%%%%%%%%%%%%%%%%%%%%%%%%%%%%%%%%%%%%%%%%%%%%%%%%%%%%%%%%%%% For fluid dynamics experiments, velocity and pressure fields are generated by numerically solving the Navier--Stokes equations. The state representation is: \begin{equation} X_t = [u(x,y,t),v(x,y,t),p(x,y,t)]. \end{equation} The computational domain is discretized into spatial elements: \begin{equation} V= \{v_1,v_2,...,v_N\}. \end{equation} Local fluid interactions define graph edges: \begin{equation} E= \{(i,j):v_i\leftrightarrow v_j\}. \end{equation} The resulting dataset contains evolving fluid graphs: \begin{equation} \mathcal{D}_{NS} = \{ G_0,G_1,...,G_T \}. \end{equation} %%%%%%%%%%%%%%%%%%%%%%%%%%%%%%%%%%%%%%%%%%%%%%%%%%%%%%%%%%%%%%%%%%%%%%%%%%%%%%%% \subsubsection{Dynamic Graph Construction} %%%%%%%%%%%%%%%%%%%%%%%%%%%%%%%%%%%%%%%%%%%%%%%%%%%%%%%%%%%%%%%%%%%%%%%%%%%%%%%% After numerical simulation, physical states are converted into graph representations. The graph construction process is: \begin{equation} X_t \rightarrow G_t=(V_t,E_t,X_t,P_t). \end{equation} Node attributes include: \begin{itemize} \item physical variables; \item local geometric information; \item material or system parameters. \end{itemize} Edge attributes include: \begin{itemize} \item distance; \item interaction strength; \item physical coupling parameters. \end{itemize} %%%%%%%%%%%%%%%%%%%%%%%%%%%%%%%%%%%%%%%%%%%%%%%%%%%%%%%%%%%%%%%%%%%%%%%%%%%%%%%% \subsubsection{Dataset Splitting} %%%%%%%%%%%%%%%%%%%%%%%%%%%%%%%%%%%%%%%%%%%%%%%%%%%%%%%%%%%%%%%%%%%%%%%%%%%%%%%% Each generated dataset is divided into training, validation, and testing subsets: \begin{equation} \mathcal{D} = \mathcal{D}_{train} \cup \mathcal{D}_{val} \cup \mathcal{D}_{test}. \end{equation} The split is performed at the trajectory level to prevent information leakage. The typical allocation is: \begin{equation} 70\%/15\%/15\% \end{equation} for training, validation, and testing. %%%%%%%%%%%%%%%%%%%%%%%%%%%%%%%%%%%%%%%%%%%%%%%%%%%%%%%%%%%%%%%%%%%%%%%%%%%%%%%% \subsubsection{Normalization and Physical Scaling} %%%%%%%%%%%%%%%%%%%%%%%%%%%%%%%%%%%%%%%%%%%%%%%%%%%%%%%%%%%%%%%%%%%%%%%%%%%%%%%% To improve numerical stability during training, physical variables are normalized: \begin{equation} \tilde{x} = \frac{x-\mu_x}{\sigma_x}, \end{equation} where $\mu_x$ and $\sigma_x$ represent the mean and standard deviation of the corresponding physical quantity. However, physical constraints are evaluated in the original dimensional space to preserve physical meaning. %%%%%%%%%%%%%%%%%%%%%%%%%%%%%%%%%%%%%%%%%%%%%%%%%%%%%%%%%%%%%%%%%%%%%%%%%%%%%%%% \subsubsection{Dataset Augmentation} %%%%%%%%%%%%%%%%%%%%%%%%%%%%%%%%%%%%%%%%%%%%%%%%%%%%%%%%%%%%%%%%%%%%%%%%%%%%%%%% To improve generalization, trajectory augmentation techniques are applied: \begin{itemize} \item variation of initial conditions; \item perturbation of physical parameters; \item temporal subsampling; \item spatial resolution changes. \end{itemize} The augmented dataset is: \begin{equation} \mathcal{D}^{*} = \mathcal{A}(\mathcal{D}), \end{equation} where $\mathcal{A}$ represents the augmentation operator. %%%%%%%%%%%%%%%%%%%%%%%%%%%%%%%%%%%%%%%%%%%%%%%%%%%%%%%%%%%%%%%%%%%%%%%%%%%%%%%% \subsubsection{Reproducibility Protocol} %%%%%%%%%%%%%%%%%%%%%%%%%%%%%%%%%%%%%%%%%%%%%%%%%%%%%%%%%%%%%%%%%%%%%%%%%%%%%%%% All dataset generation procedures are defined by: \begin{itemize} \item governing equations; \item numerical solver; \item integration timestep $\Delta t$; \item initial condition distribution; \item graph construction rules. \end{itemize} This enables complete reproduction of the experimental benchmarks. %%%%%%%%%%%%%%%%%%%%%%%%%%%%%%%%%%%%%%%%%%%%%%%%%%%%%%%%%%%%%%%%%%%%%%%%%%%%%%%% \subsubsection{Summary} %%%%%%%%%%%%%%%%%%%%%%%%%%%%%%%%%%%%%%%%%%%%%%%%%%%%%%%%%%%%%%%%%%%%%%%%%%%%%%%% The PHYSGEN experimental datasets are generated directly from physical dynamical systems and transformed into evolving graph representations. This procedure preserves the relationship between governing equations, physical interactions, and learned graph dynamics, providing a controlled environment for evaluating physics-guided long-horizon prediction. \] 
